% ============================================================
% INTEGRAZIONI CAPITOLO 2
% Integrazione framework manageriali — osservazioni Prof. Sasso
% ============================================================
%
% Questo file contiene SOLO i blocchi da inserire nel Capitolo 2.
% Non modificare il resto del capitolo.
%
% Le quattro integrazioni sono nell'ordine in cui devono apparire
% nel capitolo: tre in §2.1 (sec:ia-trasformazione) e una in §2.4
% (sec:modelli-deployment).
%
% ============================================================


% ============================================================
% INTEGRAZIONE 1 — §2.1 (prima sottosezione della sezione)
% INSERIRE dopo l'introduzione della sezione \section{IA e trasformazione
%   organizzativa} e PRIMA della prima sottosezione esistente:
%   \subsection{Un cambio di paradigma: dalla digitalizzazione
%   all'intelligenza organizzativa}
% ============================================================

\subsection{Framework interpretativi per l'adozione dell'IA nella PA}
\label{subsec:framework-interpretativi}

L'adozione dell'intelligenza artificiale nella Pubblica Amministrazione non è un fenomeno che possa essere adeguatamente compreso attraverso la sola dimensione normativa o tecnologica. Le scelte delle organizzazioni pubbliche — quali sistemi adottare, con quale modello di \textit{deployment}, in quale sequenza e con quale grado di ambizione — sono il prodotto di configurazioni complesse di fattori contestuali, dinamiche istituzionali e capacità organizzative. La letteratura manageriale e degli \textit{Information Systems} offre un insieme di framework teorici che, applicati congiuntamente, consentono di passare da una lettura descrittiva dell'adozione dell'IA a un'analisi critica e orientata all'azione. Il presente paragrafo sviluppa tre di questi framework — il \textit{TOE}, le \textit{Dynamic Capabilities} e il \textit{Strategic Alignment Model} — che saranno utilizzati come lenti interpretative nell'analisi condotta nel resto del capitolo.

\subsubsection*{Il TOE Framework applicato all'IA nella PA}

Il \textbf{\textit{Technology-Organization-Environment framework}} (\textit{TOE}), introdotto da \textcite{TornatzkyFleischer1990}, è il quadro teorico più ampiamente adottato nella letteratura sulle determinanti dell'adozione tecnologica nelle organizzazioni. La sua rilevanza per il settore pubblico è stata confermata da una sistematica \textit{review} della letteratura sull'adozione dell'IA nelle PA condotta da \textcite{MadanAshok2022} su oltre 120 studi, che individua nel \textit{TOE} il framework con maggiore potere esplicativo tra quelli applicati al contesto governativo. Anche l'analisi delle sfide sistemiche all'adozione dell'IA nelle organizzazioni sanitarie pubbliche finlandesi proposta da \textcite{PulkkinenTOE2025} mostra come le barriere all'adozione operino simultaneamente su tutti e tre i livelli del modello, rendendo insufficiente qualsiasi approccio che ne consideri uno solo.

Applicato specificamente alla scelta del modello di \textit{deployment on-premise}, il \textit{TOE} produce la seguente articolazione. Sul piano \textbf{tecnologico}, i fattori determinanti includono la disponibilità di infrastruttura hardware adeguata (server GPU, \textit{storage} ad alte prestazioni, rete interna), la compatibilità dei modelli \textit{open-source} con i sistemi informativi dell'ente, la complessità dell'implementazione e la maturità della tecnologia disponibile: la recente proliferazione di modelli di \textit{Large Language Model open-source} ad alta qualità (Llama, Mistral, DeepSeek) ha ridotto significativamente la barriera tecnologica al \textit{deployment on-premise}. Sul piano \textbf{organizzativo}, le variabili chiave sono la disponibilità di competenze interne specializzate (ingegneri ML, amministratori di sistema, figure di \textit{data governance}), la propensione della leadership alla trasformazione digitale, la capacità di \textit{change management} e la disponibilità di budget CAPEX in un contesto — quello della PA — dove i cicli di pianificazione pluriennale rendono strutturalmente difficile la conversione verso modelli di spesa OPEX puri. Sul piano \textbf{ambientale}, le pressioni più rilevanti sono di natura regolamentare: il Regolamento (UE) 2024/1689 (AI Act), la classificazione dei dati strategici prevista dalla Strategia Cloud Italia e gli obblighi della NIS~2 configurano un quadro normativo che, per determinate categorie di dati e servizi, rende la scelta \textit{on-premise} non solo preferibile ma normativamente orientata. A queste si aggiungono le pressioni istituzionali esercitate da AgID, ACN e dall'ecosistema dei fornitori di soluzioni certificate per la PA.

Lo schema seguente sintetizza l'applicazione del \textit{TOE} alla scelta \textit{on-premise}:

\begin{description}[noitemsep]
  \item[Contesto Tecnologico] Maturità LLM \textit{open-source}; disponibilità GPU; integrazione con sistemi PA; complessità implementativa.
  \item[Contesto Organizzativo] Competenze ML interne; leadership digitale; budget CAPEX; capacità di \textit{change management}.
  \item[Contesto Ambientale] AI Act; NIS~2; Strategia Cloud Italia (dati strategici); pressione AgID/ACN; best practice PSN.
\end{description}

\subsubsection*{Le Dynamic Capabilities di Teece e la PA}

Il framework delle \textbf{\textit{Dynamic Capabilities}}, elaborato da \textcite{Teece2007}, introduce una prospettiva di analisi centrata non sui fattori contestuali ma sulle \textbf{capacità organizzative} che determinano la capacità di un'organizzazione di governare attivamente — piuttosto che subire — i processi di cambiamento tecnologico. Le tre micro-fondamentali del framework — \textbf{\textit{sensing}}, \textbf{\textit{seizing}} e \textbf{\textit{transforming}} — trovano un'applicazione diretta nel contesto dell'adozione dell'IA nella PA.

La capacità di \textbf{\textit{sensing}} designa l'abilità di identificare e interpretare segnali rilevanti nell'ambiente tecnologico, regolatorio e competitivo. Per una PA, questa capacità si traduce nella facoltà di monitorare l'evoluzione del panorama dell'IA (nuovi modelli \textit{open-source}, cambiamenti normativi, \textit{best practice} europee), di riconoscere le opportunità che tale evoluzione apre per i propri processi e di anticipare i rischi regolatori prima che diventino obblighi cogenti. Come evidenziato da \textcite{RolandBergerDC2025}, le amministrazioni che hanno sviluppato capacità di \textit{sensing} avanzate sono quelle che hanno istituito unità di monitoraggio tecnologico, partecipano attivamente alle reti europee di \textit{digital government} e traducono sistematicamente gli indicatori DESI e DGI in segnali strategici.

La capacità di \textbf{\textit{seizing}} concerne la mobilitazione delle risorse per cogliere le opportunità identificate: in un contesto di PA, si traduce nella capacità di definire priorità strategiche nell'adozione dell'IA, di costruire \textit{business case} robusti per gli investimenti, di configurare processi di \textit{procurement} adeguati alla specificità delle soluzioni IA e di coinvolgere il personale nel processo di adozione. La ricerca di \textcite{DelBaroneTOE2025} sulle competenze manageriali nell'era dell'IA mostra come il \textit{seizing} richieda non solo capacità tecniche ma, soprattutto, leadership trasformativa capace di allineare gli obiettivi strategici dell'ente con le potenzialità specifiche delle soluzioni IA disponibili.

La capacità di \textbf{\textit{transforming}}, infine, riguarda la riconfigurazione continua dell'organizzazione: la revisione dei processi, la ridefinizione dei ruoli, l'aggiornamento delle competenze e l'evoluzione della cultura organizzativa. Per la PA, il \textit{transforming} è la dimensione più critica e più trascurata: l'adozione di un sistema di IA \textit{on-premise} senza una corrispondente riorganizzazione dei processi decisionali e una ridefinizione delle responsabilità degli operatori riduce l'innovazione tecnologica a un esercizio di \textit{window dressing} istituzionale.

\subsubsection*{Il Strategic Alignment Model (SAM) e la scelta infrastrutturale}

Il \textbf{\textit{Strategic Alignment Model}} (SAM), proposto da \textcite{HendersonVenkatraman1993}, offre un terzo e complementare quadro interpretativo. Il modello articola la relazione tra organizzazione e tecnologia lungo quattro domini: \textbf{\textit{business strategy}} (gli obiettivi strategici e il posizionamento dell'organizzazione nel proprio ambiente), \textbf{\textit{IT strategy}} (le scelte tecnologiche e infrastrutturali di lungo periodo), \textbf{\textit{organizational infrastructure}} (i processi, le competenze e le strutture organizzative) e \textbf{\textit{IT infrastructure}} (i sistemi, le applicazioni e i dati). Il SAM postula che la performance organizzativa sia massimizzata quando esiste un doppio allineamento: \textit{strategico} (verticale, tra \textit{business strategy} e \textit{IT strategy}) e \textit{funzionale} (orizzontale, tra i processi organizzativi e l'infrastruttura IT). Applicato al contesto della PA italiana, il SAM suggerisce che la scelta tra \textit{deployment on-premise}, \textit{cloud} o ibrido non sia una decisione puramente tecnica, ma il punto di incontro tra la strategia istituzionale dell'ente (proteggere la sovranità dei dati, garantire la continuità operativa, adempiere agli obblighi NIS~2) e la sua strategia IT (qualificazione dei sistemi, interoperabilità, ciclo di vita dei modelli IA). Come dimostra \textcite{CordeiroSAM2019}, l'applicazione del SAM al settore pubblico rivela come le scelte infrastrutturali disallineate rispetto alla strategia istituzionale dell'ente producano sistematicamente inefficienze, duplicazioni e rischi di \textit{lock-in} tecnologico. La dimensione dell'allineamento strategico sarà sviluppata in modo specifico nella sezione~\ref{sec:modelli-deployment}, dove si analizzeranno le implicazioni del SAM per la scelta del modello di \textit{deployment} dell'IA.


% ============================================================
% INTEGRAZIONE 2 — §2.1 (seconda nuova sottosezione)
% INSERIRE dopo le sottosezioni esistenti:
%   - "Un cambio di paradigma..."
%   - "Dall'automazione al supporto decisionale strategico"
%   - "Rischi organizzativi e resistenza al cambiamento"
%   - "Il Piano Triennale e gli obiettivi quantitativi per l'IA nella PA"
% e PRIMA della riga:
%   "In sintesi, l'IA si configura come un fattore di trasformazione..."
% (oppure inserire come ultima sottosezione di §2.1, subito prima
%  della transizione verso §2.2)
% ============================================================

\subsection{IA come leva strategica e decisionale}
\label{subsec:ia-leva-strategica}

L'analisi condotta nei paragrafi precedenti ha evidenziato come l'IA nella PA sia stata per lungo tempo confinata al ruolo di strumento di automazione operativa. Una lettura più critica della letteratura recente suggerisce che questa prospettiva sia non soltanto riduttiva, ma strategicamente fuorviante: l'intelligenza artificiale rappresenta anzitutto una \textbf{leva di trasformazione del processo decisionale strategico}, con implicazioni che investono la struttura stessa del governo pubblico.

\textcite{ShresthaAISDM2024}, in un'analisi pubblicata su \textit{Strategy Science}, mostrano come l'IA possa ampliare le capacità decisionali strategiche lungo tre dimensioni fondamentali. La prima è la \textbf{\textit{search augmentation}}: i sistemi di IA sono in grado di esplorare spazi di soluzione vastissimi, identificando opzioni che sfuggirebbero all'analisi umana tradizionale, e possono simulare virtualmente scenari strategici complessi prima che vengano implementati. La seconda dimensione è la \textbf{\textit{representation augmentation}}: l'IA consente di costruire rappresentazioni più ricche e multidimensionali dei problemi strategici, integrando fonti eterogenee di dati (strutturati e non strutturati, interni ed esterni) in modelli interpretativi coerenti. La terza è la \textbf{\textit{aggregation augmentation}}: i sistemi di IA possono sintetizzare e ponderare le valutazioni di attori multipli, riducendo i bias cognitivi individuali e migliorando la qualità delle decisioni collettive. Applicata al contesto della PA, questa triplice capacità di augmentazione strategica si traduce nella possibilità per le amministrazioni di affinare la pianificazione dei servizi, la previsione dei fabbisogni e la valutazione dell'impatto delle politiche pubbliche con un grado di sofisticazione analitica precedentemente accessibile solo a organizzazioni private dotate di cospicue risorse di \textit{data science}.

Una seconda prospettiva utile a comprendere la dimensione strategica dell'IA nella PA è offerta dal \textbf{framework macro-meso-micro} elaborato da \textcite{ValleCruzMacroMesoMicro2024}. Gli autori propongono di analizzare l'impatto dell'IA nel settore pubblico su tre livelli distinti ma interdipendenti. Al livello \textbf{macro}, l'IA ridisegna i regimi di governance, le strategie nazionali e i sistemi di regolazione: l'AI Act europeo, la Strategia Italiana per l'IA e i target del Decennio Digitale 2030 sono espressione di questo livello. Al livello \textbf{meso}, l'IA trasforma le organizzazioni pubbliche nella loro struttura interna, i rapporti tra enti, i network settoriali e le catene del valore dei servizi pubblici: è a questo livello che si collocano le scelte di \textit{deployment}, i modelli di \textit{procurement} e le architetture di interoperabilità. Al livello \textbf{micro}, l'IA modifica i comportamenti dei \textit{street-level bureaucrats}, dei dirigenti pubblici e dei cittadini, ridefinendo i confini della discrezionalità amministrativa e le modalità di interazione con i servizi pubblici. Il valore analitico di questo framework risiede nel mettere in luce come le decisioni adottate a livello meso — inclusa la scelta del modello di \textit{deployment on-premise} — non siano neutrali rispetto ai processi in corso ai livelli macro (conformità al quadro regolatorio europeo) e micro (qualità dell'esperienza operativa del personale e del cittadino).

Questa lettura multi-livello consente di ricuperare e approfondire il concetto di \textbf{\textit{proactiveness}} introdotto dal \textit{Digital Government Policy Framework} OCSE nel Capitolo~1: la PA proattiva non è semplicemente quella che «usa l'IA» per automatizzare le pratiche, ma quella che ha sviluppato la capacità — a livello macro, meso e micro — di impiegare l'intelligenza artificiale come strumento di anticipazione strategica, di ridisegno dei processi e di personalizzazione dei servizi. Sotto questo profilo, la scelta del modello infrastrutturale — e in particolare la capacità di mantenere il controllo sui modelli, sui dati e sull'inferenza in ambiente \textit{on-premise} — è una precondizione per esercitare pienamente questa \textit{proactiveness}, perché garantisce l'indipendenza operativa necessaria ad adattare continuamente i sistemi di IA alle esigenze specifiche dell'ente senza dipendere dai vincoli imposti da fornitori terzi.


% ============================================================
% INTEGRAZIONE 3 — §2.1 (terza nuova sottosezione)
% INSERIRE subito DOPO la sottosezione "IA come leva strategica e decisionale"
% (e prima della transizione verso §2.2)
% ============================================================

\subsection{Pressioni istituzionali e isomorfismo nell'adozione dell'IA}
\label{subsec:isomorfismo-ia}

La comprensione delle ragioni per cui le pubbliche amministrazioni adottano — o resistono all'adozione de — l'intelligenza artificiale non può limitarsi a considerazioni di efficienza o di conformità normativa. La \textbf{\textit{Institutional Theory}}, e in particolare il concetto di \textbf{isomorfismo istituzionale} elaborato da \textcite{DiMaggioPow1983}, offre una prospettiva analitica complementare che illumina le dinamiche di imitazione, pressione e professionalizzazione che operano nel «campo organizzativo» della PA.

Il primo meccanismo — l'isomorfismo \textbf{coercitivo} — è il più visibile e il più direttamente misurabile. Le PA italiane ed europee sono sottoposte a un sistema di obblighi normativi crescentemente denso in materia di IA: il Regolamento (UE) 2024/1689 (AI Act) definisce categorie di rischio e obblighi di conformità che rendono l'adozione dell'IA un processo governato ab initio da requisiti cogenti \parencite{AIAct2024}; la Direttiva NIS~2, recepita con il D.Lgs. 138/2024, impone standard di sicurezza informatica che condizionano l'architettura dei sistemi; la Strategia Cloud Italia vincola le scelte di \textit{deployment} dei dati strategici a infrastrutture qualificate nazionali. In questo quadro, la pressione coercitiva non spinge soltanto verso l'adozione dell'IA, ma ne orienta le modalità: la conformità normativa diventa una forza di selezione che favorisce determinati modelli di \textit{deployment} — in particolare quelli che garantiscono piena tracciabilità, controllo dei dati e \textit{auditability} — rispetto ad altri.

Il secondo meccanismo — l'isomorfismo \textbf{mimetico} — opera nelle condizioni di incertezza tecnologica e organizzativa che caratterizzano l'adozione di tecnologie emergenti. In assenza di pratiche consolidate, le PA tendono a imitare le scelte degli enti percepiti come più innovativi o più autorevoli: le amministrazioni «follower» replicano le architetture del Polo Strategico Nazionale, adottano le soluzioni pilota diffuse da AgID o si conformano alle esperienze degli enti «campioni» individuati dalla Commissione Europea nell'\textit{eGovernment Benchmark}. Come documentato da \textcite{PatalonIsomorphism2024} nell'analisi dei processi di trasformazione digitale nei comuni europei, il mimetismo istituzionale può accelerare la diffusione delle innovazioni, ma tende anche a produrre omologazione delle soluzioni adottate anche quando le condizioni contestuali degli enti sarebbero diverse, generando disallineamenti tra architettura scelta e reale configurazione \textit{TOE} dell'organizzazione.

Il terzo meccanismo — l'isomorfismo \textbf{normativo} — è prodotto dai processi di professionalizzazione e dalla diffusione di standard condivisi attraverso reti di esperti e percorsi formativi strutturati. Il ruolo di AgID Academy, i programmi CRUI per la formazione dei dirigenti pubblici sull'IA, le linee guida operative del Piano Triennale 2024--2026 e i percorsi di certificazione ISO/IEC 42001 contribuiscono a definire, nel campo della PA digitale, un insieme di «pratiche professionali legittime» che orientano le scelte tecnologiche verso modelli conformi alle aspettative della comunità professionale di riferimento. La tensione tra questi tre meccanismi isomorfi — normatività regolatoria, imitazione dei \textit{best performer} e professionalizzazione — spiega alcune apparenti contraddizioni nell'adozione dell'IA nella PA: la coesistenza di enti altamente innovativi e enti ancora fermi alle fasi sperimentali, all'interno di un medesimo quadro normativo nazionale, riflette differenti esposizioni ai tre meccanismi e differenti capacità di elaborarne le pressioni.


% ============================================================
% INTEGRAZIONE 4 — §2.4
% INSERIRE dopo \subsection{Il modello ibrido a tre livelli}
% (che termina con: "...e ricorrendo al cloud per i picchi temporanei
%  e per le fasi di sperimentazione.")
% e PRIMA di \subsection{Verso il capitolo successivo}
% oppure, se non esiste questa sottosezione, PRIMA del paragrafo conclusivo
% che inizia con: "In conclusione, la scelta del modello di deployment..."
% ============================================================

\subsection{Allineamento strategico e scelta del modello di deployment}
\label{subsec:allineamento-strategico}

L'analisi dei modelli di \textit{deployment} — condotta nei paragrafi precedenti lungo le dimensioni del costo, della sicurezza, della latenza e della conformità normativa — rischia di essere interpretata in modo riduttivo se non viene inquadrata in una prospettiva di \textbf{allineamento strategico}. La scelta tra \textit{on-premise}, \textit{cloud} e \textit{hybrid} non è, in ultima analisi, una decisione tecnica delegabile al reparto IT: è una \textbf{decisione strategica} che richiede un esplicito allineamento tra la strategia istituzionale della PA e la sua strategia tecnologica.

Il \textbf{\textit{Strategic Alignment Model}} di \textcite{HendersonVenkatraman1993}, introdotto nel paragrafo~\ref{subsec:framework-interpretativi}, fornisce il quadro analitico per comprendere perché questo allineamento sia sia necessario sia strutturalmente difficile da realizzare nella PA. Il modello distingue, sul versante dell'organizzazione, tra \textit{business strategy} (gli obiettivi istituzionali, il posizionamento nella rete delle PA e le priorità di servizio) e \textit{organizational infrastructure} (i processi, le competenze e le strutture); sul versante tecnologico, tra \textit{IT strategy} (le scelte di lungo periodo in materia di piattaforme, architetture e fornitori) e \textit{IT infrastructure} (i sistemi operativi, le applicazioni e i dati). Il SAM postula che la performance sia massimizzata quando i quattro domini sono mutuamente allineati, tanto in senso verticale (strategia istituzionale con strategia IT) quanto in senso orizzontale (processi organizzativi con infrastruttura tecnica). Come dimostra \textcite{CordeiroSAM2019}, applicando il SAM al settore pubblico, le PA che hanno prodotto i risultati più significativi nella transizione digitale sono quelle che hanno trattato le scelte infrastrutturali come estensioni della propria strategia istituzionale, non come variabili operative autonome.

Applicato alla scelta del modello di \textit{deployment} dell'IA, il SAM produce un'indicazione analitica precisa. Una PA che ha definito, a livello di \textit{business strategy}, obiettivi di \textbf{sovranità sui dati}, \textbf{autonomia operativa} e \textbf{conformità agli obblighi NIS~2} deve necessariamente riflettere questi obiettivi nella propria \textit{IT strategy}: la scelta di un \textit{deployment cloud} presso un operatore extra-UE genererebbe un disallineamento strutturale che nessuna misura di sicurezza compensativa potrebbe eliminare del tutto, dato il regime giurisdizionale del \textit{CLOUD Act}. Al contrario, un'architettura \textit{on-premise} o ibrida con livello critico \textit{on-premise} produce un allineamento coerente tra la strategia istituzionale (sovranità, continuità operativa, \textit{auditability}) e la strategia tecnologica (\textit{deployment} controllato, \textit{logging} interno, piena accessibilità ai modelli). Sul piano orizzontale, l'allineamento funzionale richiede che l'infrastruttura IT scelta sia coerente con i processi organizzativi dell'ente: un sistema di IA \textit{on-premise} di alto profilo computazionale risulta disallineato se l'organizzazione non dispone delle competenze interne necessarie per gestirlo e aggiornarlo nel tempo, generando dipendenza da fornitori di \textit{managed service} che possono ricostituire, sotto altra forma, il rischio di \textit{lock-in} che si intendeva evitare.

La prospettiva del SAM consente dunque di superare la dicotomia tecnica tra \textit{on-premise} e \textit{cloud} e di ricondurla a una questione di \textbf{coerenza strategica}: il modello di \textit{deployment} ottimale per una PA non è quello che massimizza una singola dimensione (costo, sicurezza, scalabilità) ma quello che produce il maggiore allineamento tra la strategia istituzionale dell'ente, la sua strategia IT e le sue capacità organizzative effettive. Il modello ibrido a tre livelli analizzato nel paragrafo precedente può essere letto, in questa prospettiva, come la configurazione architetturale che offre il maggiore spazio per un allineamento adattivo: la segmentazione per criticità dei dati consente di allineare ciascun livello dell'architettura al corrispondente profilo di rischio e agli obiettivi strategici dell'ente, senza imporre una soluzione unica che risulterebbe inevitabilmente disallineata rispetto a qualcuna delle dimensioni rilevanti.
